\documentclass[10pt,letterpaper]{article}


\usepackage[margin=0.5in]{geometry}
\usepackage{fontspec}
\setmainfont{Latin Modern Roman}
\usepackage{amsmath,amssymb,mathtools}
\usepackage{booktabs,array,longtable}
\usepackage{enumitem}
\usepackage{xcolor}
\usepackage{fancyhdr}
\usepackage[hidelinks]{hyperref}
\usepackage{microtype}

\definecolor{courseblue}{RGB}{41,77,115}
\definecolor{coursegray}{RGB}{92,105,117}

\renewcommand{\arraystretch}{1.15}


\newcommand{\questionheading}[1]{%
  \subsection*{Question #1}
}
\newcommand{\solutionheading}[2]{%
  \subsection*{Solution #1: #2}
}
\newenvironment{answer}{%
  \begin{quote}\small\color{black}
}{%
  \end{quote}
}

\title{\textbf{STAT 1210-2}\\[2mm]
\large Section 2: Probability and Distributions Exercises}
\date{}
\begin{document}
\maketitle
\questionheading{1}
A network of wildlife cameras operated for 1200 camera-nights. A red fox was recorded during 174 of those nights.
\begin{enumerate}
	\item Use the observations to estimate the probability that a red fox is recorded during a randomly selected camera-night.
	\item Estimate the probability that a red fox is not recorded during a randomly selected camera-night.
	\item If similar conditions apply during another 500 camera-nights, how many nights would be expected to contain a red-fox detection?
	\item Does the answer to Part 3 guarantee the observed number of detections? Explain using the frequentist interpretation of probability.
\end{enumerate}

\questionheading{2}
Two laboratory assays are performed on the same specimen. Each assay is recorded as positive $(+)$ or negative $(-)$, and the order of the assays is retained.
\begin{enumerate}
	\item Write the sample space for the two assay results.
	\item Write the event that exactly one assay is positive.
	\item Write the event that at least one assay is positive.
	\item Write the complement of the event in Part 3.
	\item If each assay has probability $0.20$ of a positive result and the assay results are independent, calculate the probability that exactly one assay is positive.
\end{enumerate}

\questionheading{3}
Let $X$ be the number of seedlings, out of four, that show measurable root growth after a laboratory treatment. A proposed probability model is shown below.
\begin{center}
\begin{tabular}{c*{5}{r}}
\toprule
$x$ & 0 & 1 & 2 & 3 & 4 \\
\midrule
$P(X=x)$ & 0.10 & 0.24 & 0.32 & 0.22 & $c$ \\
\bottomrule
\end{tabular}
\end{center}
\begin{enumerate}
	\item Find $c$ so that the table is a valid probability model.
	\item Calculate $P(X\geq 3)$ and $P(X<2)$.
	\item Calculate the mean $\mu_X=E(X) = \sum xP(X=x)$.
	\item Calculate the variance and standard deviation of $X$. Note that you can also calculate variance using this formula
	$$\operatorname{Var}(X)=E(X^2)-[E(X)]^2 = \sum x^2P(X=x) - \left(\sum xP(X=x) \right)^2$$
	\item Interpret the mean in context. Explain why the mean does not have to be a possible observed count.
\end{enumerate}

\newpage 

\questionheading{4}
The calibration error $Y$ of a timing sensor is uniformly distributed between $-0.6$ and $0.6$ microseconds.
\begin{enumerate}
	\item Find the height of the uniform density curve.
	\item Use areas to calculate $P(-0.2\leq Y\leq 0.3)$.
	\item Use areas to calculate $P(Y<-0.45\text{ or }Y>0.50)$.
	\item Find $P(Y=0)$ and explain the result.
	\item Explain why $P(Y\leq0.3)=P(Y<0.3)$.
	\end{enumerate}

\questionheading{5}
An insurer uses the following simplified probability model for the annual claim cost $C$ of one policy.
\begin{center}
\begin{tabular}{c*{4}{r}}
\toprule
Claim cost $c$ & \$0 & \$500 & \$3000 & \$15,000 \\
\midrule
$P(C=c)$ & 0.88 & 0.08 & 0.03 & 0.01 \\
\bottomrule
\end{tabular}
\end{center}
\begin{enumerate}
	\item Calculate the expected annual claim cost $E(C)$.
	\item Calculate $E(C^2)$.
	\item Use $\operatorname{Var}(C)=E(C^2)-[E(C)]^2$ to calculate the variance and standard deviation.
	\item Find the expected total claim cost for 10,000 policies under this model.
	\item If the insurer collects a premium of \$350 per policy before administrative expenses, calculate the expected amount remaining per policy after claim costs. Explain why this is a long-run average rather than a guaranteed profit on every policy.
\end{enumerate}

\questionheading{6}
Among 2500 animals receiving a veterinary medication, 37 experienced a specified adverse reaction.
\begin{enumerate}
	\item Calculate the empirical risk of the adverse reaction.
	\item Calculate the empirical probability of no adverse reaction.
	\item Calculate the odds of an adverse reaction and express them in the form \`\`about 1 to $k$.''
	\item Explain the difference between a risk of $0.0148$ and odds of approximately $1:66.6$.
\end{enumerate}

\questionheading{7}
Let $X$ be the number of surviving offspring in a four-offspring litter. The probability model is
\begin{center}
\begin{tabular}{c*{5}{r}}
\toprule
$x$ & 0 & 1 & 2 & 3 & 4 \\
\midrule
$P(X=x)$ & 0.05 & 0.15 & 0.30 & 0.30 & 0.20 \\
\bottomrule
\end{tabular}
\end{center}
A rehabilitation program models the litter's cost, in dollars, as
\[
	Y=120+45X
\]
\begin{enumerate}
	\item Calculate $E(X)$ and $\operatorname{SD}(X)$ (Note: SD here means standard deviation).
	\item Calculate $E(Y)$ and $\operatorname{SD}(Y)$ without constructing a new probability table.
	\item Calculate $P(Y\geq255)$.
	\item Interpret $E(Y)$ in context.
\end{enumerate}

\newpage

\questionheading{8}
In a population of dogs, let $V$ be the event that a dog is vaccinated and let $I$ be the event that the dog carries a particular infection. Express each statement using set notation.
\begin{enumerate}
	\item The dog is vaccinated and infected.
	 \item The dog is vaccinated or infected, including the possibility of both.
	\item The dog is vaccinated but not infected.
	\item The dog is neither vaccinated nor infected.
	\item The dog has at least one of the two characteristics.
\end{enumerate}

\questionheading{9}
Researchers examine 400 seabirds. Of these birds, 145 show evidence of plastic ingestion, 92 have low body mass and 55 have both characteristics. Let $A$ denote plastic ingestion and $B$ denote low body mass.
\begin{enumerate}
	\item Calculate $P(A)$, $P(B)$ and $P(A\cap B)$.
	\item Use the general addition rule to calculate $P(A\cup B)$.
	\item Calculate the probability that a bird has neither characteristic.
	\item Calculate the probability that a bird has plastic ingestion but not low body mass.
	 \item Explain why simply adding $P(A)$ and $P(B)$ would give an incorrect answer for $P(A\cup B)$.
\end{enumerate}

\questionheading{10}
In a population, the probability that an individual carries genetic marker $A$ is $0.08$, and the probability of carrying marker $B$ is $0.12$. Assume the two marker events are independent.
\begin{enumerate}
	\item Calculate the probability that an individual carries both markers.
	\item Calculate the probability that an individual carries at least one marker.
	\item Calculate the probability that an individual carries neither marker.
	\item Are the events \`\`carries marker $A$'' and \`\`carries marker $B$'' disjoint? Explain.
\end{enumerate}

\questionheading{11}
A veterinary hospital records the following results for 300 dogs.
\begin{center}
\begin{tabular}{lrrr}
\toprule
	& Diarrhea & No diarrhea & Total \\
\midrule
Parasitic infection & 48 & 32 & 80 \\
No parasitic infection & 36 & 184 & 220 \\
\midrule
Total & 84 & 216 & 300 \\
\bottomrule
\end{tabular}
\end{center}

Let $I$ denote parasitic infection and $D$ denote diarrhea.
\begin{enumerate}
	\item Calculate $P(I\cap D)$.	
	\item Calculate and interpret $P(I\mid D)$.
	\item Calculate and interpret $P(D\mid I)$.
	\item Explain why $P(I\mid D)$ and $P(D\mid I)$ are different even though both use the same intersection count.
\end{enumerate}

\newpage

\questionheading{12}
A wildlife health study records environmental exposure and skin lesions in 300 animals.
\begin{center}
\begin{tabular}{lrrr}
\toprule
	& Lesion & No lesion & Total \\
\midrule
Exposed & 30 & 70 & 100 \\
Not exposed & 20 & 180 & 200 \\
\midrule
Total & 50 & 250 & 300 \\
\bottomrule
\end{tabular}
\end{center}
Let $E$ denote exposure and $L$ denote a lesion.
\begin{enumerate}
	\item Calculate $P(L)$ and $P(L\mid E)$.
	\item Based on these probabilities, are $E$ and $L$ independent? Explain.
	\item Check independence a second way by comparing $P(E\cap L)$ with $P(E)P(L)$.
	\item Does the observed association prove that the exposure caused the lesions? Explain briefly.
\end{enumerate}

\questionheading{13}
In a population of plants, 4\% carry a pathogen. A screening test has sensitivity 92\% and specificity 95\%. Consider 10,000 plants representative of this population.
\begin{enumerate}
	\item Construct the expected counts of true positives, false negatives, false positives and true negatives.
	\item Calculate the overall probability of a positive result.
	\item Calculate $P(\text{pathogen}\mid +)$ and interpret it.
	\item Calculate $P(\text{pathogen}\mid -)$ and interpret it.
	\item Explain why high sensitivity and specificity do not necessarily produce a positive predictive value near 100\% when prevalence is low.
\end{enumerate}

\questionheading{14}
Suppose a DNA profile would match a randomly selected innocent person with probability $1/10{,}000$. A city contains approximately 100,000 people who could not initially be excluded as possible suspects.
\begin{enumerate}
	\item Among 100,000 innocent people, how many coincidental matches would be expected under this simplified model?	
	\item Explain the difference between
	\[
	P(\text{match}\mid\text{innocent}) \quad\text{and}\quad P(\text{innocent}\mid\text{match})
	\]
	\item Can the second probability be calculated from the match probability alone? Explain what additional information is needed.
	\item State the error made in the prosecutor's fallacy.
\end{enumerate}

\questionheading{15}
An animal shelter has 12 cats, of which 8 are vaccinated and 4 are not vaccinated. Three different cats are selected at random without replacement.
\begin{enumerate}
	\item Calculate the probability that all three selected cats are vaccinated.
	\item Calculate the probability that at least one selected cat is not vaccinated.
	 \item Calculate the probability that exactly two selected cats are vaccinated.
	\item Given that the first two selected cats are vaccinated, calculate the probability that the third is vaccinated.
	\item Explain why the three selection outcomes are not independent.
\end{enumerate}

\newpage

\questionheading{16}
Resting heart rates in a population of healthy adult horses are approximately Normal with mean 40 beats per minute and standard deviation 6 beats per minute.
\begin{enumerate}
\item Find the $z$-score for a horse with a resting heart rate of 52 beats per minute.
\item Find the $z$-score for a horse with a resting heart rate of 34 beats per minute.
\item Interpret both $z$-scores in context.
\item Which observation is more unusual relative to the population distribution? Explain.
\end{enumerate}

\questionheading{17}
Let $Z$ have the standard Normal distribution. Use a standard Normal table to calculate each probability.
\begin{enumerate}
\item $P(-1.25<Z<0.80)$
\item $P(Z>1.65)$
\item $P(|Z|>2.00)$
\item Sketch a standard Normal curve and shade the region corresponding to each probability.
\end{enumerate}

\questionheading{18}
The fork length $X$ of adult Atlantic salmon returning to a river is approximately Normally distributed with mean 50 cm and standard deviation 6 cm.
\begin{enumerate}
\item Calculate $P(X<44)$.
\item Calculate $P(X>59)$.
\item Calculate $P(47<X<56)$.
\item In a group of 800 returning salmon, approximately how many would be expected to exceed 59 cm?
\item Explain why the expected count in Part 4 need not equal the observed count exactly.
\end{enumerate}

\questionheading{19}
An enzyme measurement in a healthy population is approximately Normal with mean 120 units and standard deviation 15 units.
\begin{enumerate}
\item Find the 90th percentile of the distribution.
\item Find the value that separates the lowest 5\% of measurements from the remaining 95\%.
\item Find the two values that contain the central 95\% of measurements.
\item Interpret the 90th percentile in context.
\end{enumerate}

\questionheading{20}
A dog's fasting glucose measurement is 135 mg/dL. In the relevant healthy dog population, glucose has mean 100 mg/dL and standard deviation 15 mg/dL. A cat's creatinine measurement is 2.1 mg/dL. In the relevant healthy cat population, creatinine has mean 1.4 mg/dL and standard deviation 0.35 mg/dL. Assume both reference distributions are approximately Normal.
\begin{enumerate}
\item Calculate the $z$-score for each measurement.
\item Which measurement is farther above its population mean in standard-deviation units?
\item Calculate the upper-tail probability associated with each value.
\item Explain why the raw measurements cannot be compared directly.
\end{enumerate}

\newpage

\questionheading{21}
The wing span of a bird species is approximately Normal with mean 82 cm and standard deviation 4 cm.
\begin{enumerate}
	\item Use the 68--95--99.7 rule to give the interval containing approximately 68\% of wing spans.
	\item Give the interval containing approximately 95\% of wing spans.
	\item Give the interval containing approximately 99.7\% of wing spans.
	\item Among 2000 birds, approximately how many would be expected to have wing spans outside the 95\% interval?
	\item Approximately how many would be expected to have wing spans greater than 90 cm?
\end{enumerate}

\questionheading{22}
Recovery time $X$ after a veterinary procedure is approximately Normal with mean 50 minutes and standard deviation 8 minutes.
\begin{enumerate}
	\item Calculate $P(X>60)$.
	\item Calculate $P(X>45)$.
	\item Among animals whose recovery time exceeds 45 minutes, calculate the probability that recovery exceeds 60 minutes.
	\item Find the recovery time that marks the slowest 10\% of recoveries.
	\item Explain the difference between the unconditional probability in Part 1 and the conditional probability in Part 3.
\end{enumerate}

\questionheading{23}
For each setting, state whether the count is reasonably modelled as binomial, Poisson or neither. Give a brief reason.
\begin{enumerate}
	\item The number of 20 independently treated patients who experience an adverse reaction, when each patient has the same reaction probability.
	\item The number of emergency calls received during one hour, assuming a stable average rate and independent occurrences.
	\item The number of tagged fish in a sample of 10 fish drawn without replacement from a small tank containing 15 tagged and 15 untagged fish.
	\item The number of patients treated until the first successful response occurs.
	\item The number of bacterial colonies appearing on a fixed-area plate, assuming colonies arise independently at a stable average rate across the plate.
\end{enumerate}

\questionheading{24}
Each of 12 independently treated animals has probability 0.80 of producing a protective immune response. Let $X$ be the number producing a protective response.
\begin{enumerate}
	\item State the probability distribution of $X$.
	\item Calculate $P(X=10)$ using the binomial formula.
	\item Calculate the probability that all 12 animals produce a protective response.
	\item Find the mean and standard deviation of $X$.
	\item Interpret the mean in context.
\end{enumerate}

\newpage

\questionheading{25}
A manufacturer reports that each diagnostic cartridge has probability 0.04 of being defective. Assume cartridge outcomes are independent. A laboratory receives 20 cartridges, and $X$ is the number that are defective.
\begin{enumerate}
	\item State the binomial model for $X$.
	\item Calculate $P(X=0)$ and $P(X=1)$.
	\item Calculate $P(X\geq2)$ using a complement.
	\item Find the mean and standard deviation of $X$.
	\item Explain why the mean does not imply that a box can contain 0.8 of a defective cartridge.
\end{enumerate}

\questionheading{26}
The number $X$ of policyholders making a specified type of claim in a group of 50 independent policyholders is modelled as binomial. The expected number of claims is 5.
\begin{enumerate}
	\item Use the binomial mean to find the claim probability $p$.
	\item Find the standard deviation of $X$.
	\item Calculate the probability of no claims.
		\item Calculate the probability of at most two claims.
	\item Explain what assumption would be questionable if a single storm could cause many of the policyholders to claim at the same time.
\end{enumerate}

\questionheading{27}
Wildlife--vehicle collisions on a rural highway are modelled by a Poisson process with an average of 3.6 collisions per week. Let $X$ be the number of collisions during a two-day period.
\begin{enumerate}
	\item Find the appropriate value of $\lambda$ for two days.
	\item Calculate $P(X=0)$.
	\item Calculate $P(X=1)$.
	\item Calculate $P(X\geq2)$.
	\item State the mean and standard deviation of the two-day count.
\end{enumerate}

\questionheading{28}
Under a baseline model, the monthly number $X$ of seabird deaths found along a monitored coastline follows a Poisson distribution with $\lambda=2$.
\begin{enumerate}
	\item Calculate $P(X\leq1)$.
	\item Calculate $P(X\geq7)$.
	\item Find the smallest whole number $k$ for which $P(X\geq k)<0.01$.
	\item If seven deaths are observed, explain what the small probability does and does not establish.
\end{enumerate}

\newpage

\questionheading{29}
The number of ticks collected along a 100-metre forest transect is modelled as Poisson with mean 1.5. Assume the same rate applies along a comparable 300-metre transect.
\begin{enumerate}
	\item Find the Poisson parameter for the 300-metre transect.
	\item Calculate the probability of exactly five ticks along 300 metres.
	\item Calculate the probability of at least eight ticks along 300 metres.
	\item Find the mean, variance and standard deviation of the 300-metre count.
	\item Identify one real-world feature that could make the Poisson model inappropriate.
\end{enumerate}

\questionheading{30}
In a population of animals, 12\% carry an infection. A screening test has sensitivity 90\% and specificity 95\%. Fifteen animals are independently selected from the population and tested. Let $Y$ be the number of positive test results.
\begin{enumerate}
	\item Use the law of total probability to find the probability that one randomly selected animal tests positive.
	\item State an appropriate probability model for $Y$.
	\item Calculate $P(Y\geq3)$.
	\item Find the mean and standard deviation of $Y$.
	\item Given that one randomly selected animal tests positive, calculate the probability that it actually carries the infection.
	\item Explain why the probability in Part 5 differs from the sensitivity.
\end{enumerate}

\questionheading{31}
A marine research team tags Atlantic salmon before migration. Let $R$ be the event that a tagged salmon returns to the river the following year, and let $T$ be the event that its transmitter is still functioning when it returns. Suppose
\[
	P(R)=0.62, \qquad P(T)=0.80, \qquad P(R\cap T)=0.54
\]
\begin{enumerate}
	\item Calculate $P(R\cup T)$.
	\item Calculate $P(R\mid T)$ and interpret the result in context.
	\item Calculate the probability that neither event occurs.
	\item Are $R$ and $T$ disjoint? Explain.
	\item Are $R$ and $T$ independent? Support your answer with a calculation.
\end{enumerate}

\questionheading{32}
Let $X$ be the number of emergency patients arriving at a veterinary clinic during a randomly selected two-hour period. The proposed probability model is shown below.

\begin{center}
\begin{tabular}{c*{5}{r}}
\toprule
$x$ & 0 & 1 & 2 & 3 & 4 \\
\midrule
$P(X=x)$ & 0.10 & 0.25 & 0.35 & 0.20 & $c$ \\
\bottomrule
\end{tabular}
\end{center}

\begin{enumerate}
\item Find $c$ so that this is a valid probability model.
\item Calculate $P(X\geq2)$ and $P(1\leq X<4)$.
\item Calculate the mean $E(X)$.
\item Calculate the variance and standard deviation of $X$.
\item Interpret $E(X)$ in the context of many two-hour periods.
\end{enumerate}

\newpage

\questionheading{33}
A current meter used in an oceanography experiment has a measurement error $Y$, in centimetres per second, that is uniformly distributed between $-0.3$ and $0.5$.
\begin{enumerate}
	\item Find the height of the uniform density curve.
	\item Calculate $P(-0.1\leq Y\leq0.2)$ using the area under the density curve.
	\item Calculate $P(Y<-0.2\text{ or }Y>0.35)$.
	\item Calculate $P(Y<0.1\mid Y>-0.1)$.
	\item What is $P(Y=0.1)$? Explain using the geometry of a continuous probability distribution.
\end{enumerate}

\questionheading{34}
An insurance company uses the following simplified model for a pet owner's eligible veterinary expenses $C$ during one year.
\begin{center}
\begin{tabular}{c*{4}{r}}
\toprule
Eligible expenses $c$ & \$0 & \$400 & \$1200 & \$4000 \\
\midrule
$P(C=c)$ & 0.55 & 0.25 & 0.15 & 0.05 \\
\bottomrule
\end{tabular}
\end{center}

The policy has a \$500 annual deductible and pays all eligible expenses above the deductible. The insurer's payment is therefore $Y=\max(C-500,0)$.
\begin{enumerate}
	\item Construct the probability distribution of $Y$.
	\item Calculate the expected annual eligible expense $E(C)$.
	\item Calculate the insurer's expected annual payment $E(Y)$.
	\item Calculate the owner's expected annual out-of-pocket amount $E(C-Y)$.
	\item Explain why the expected payment is not the amount the insurer will pay for every policy.
\end{enumerate}

\questionheading{35}
A wildlife health study records whether each of 500 seabirds had substantial exposure to an environmental contaminant and whether a particular liver lesion was present.
\begin{center}
\begin{tabular}{lrrr}
\toprule
 & Lesion present & Lesion absent & Total \\
\midrule
Substantial exposure & 72 & 108 & 180 \\
No substantial exposure & 48 & 272 & 320 \\
\midrule
Total & 120 & 380 & 500 \\
\bottomrule
\end{tabular}
\end{center}

Let $E$ denote substantial exposure and let $L$ denote the presence of the lesion.
\begin{enumerate}
	\item Calculate $P(L\mid E)$ and $P(L\mid E^c)$.
	\item Calculate the risk ratio $P(L\mid E)/P(L\mid E^c)$ and interpret it.
	\item Calculate $P(E\mid L)$.
	\item Are $E$ and $L$ independent? Support your answer numerically.
	\item Can this observational study establish that the contaminant caused the lesion? Explain one limitation.
\end{enumerate}

\newpage

\questionheading{36}
In a population of cats, 8\% have early kidney disease. A screening test has sensitivity 91\% and specificity 94\%.
\begin{enumerate}
	\item For a hypothetical population of 10,000 cats, find the expected numbers of true positives, false negatives, false positives and true negatives.
	\item Calculate the probability that a randomly selected cat tests positive.
	\item Given a positive result, calculate the probability that the cat has early kidney disease.
	\item Given a negative result, calculate the probability that the cat does not have early kidney disease.
	\item Explain why false positives may outnumber false negatives even though the test has high sensitivity and specificity.
\end{enumerate}

\questionheading{37}
Daily milk production for cows in a particular herd is approximately Normally distributed with mean 31.5 litres and standard deviation 4.2 litres.
\begin{enumerate}
	\item Calculate the probability that a randomly selected cow produces between 28 and 36 litres in one day.
	\item Find the daily production level that separates the highest 10\% of cows from the remaining 90\%.
	\item Among 800 cows, approximately how many would be expected to produce less than 24 litres in one day?
	\item If 2.5 litres are added to every recorded value because of a calibration correction, state the new mean and standard deviation.
	\item Explain why the calculation in Part 4 changes the mean but not the standard deviation.
\end{enumerate}

\questionheading{38}
A wildlife vaccination program distributes oral vaccine baits. Under similar field conditions, each of 14 independently monitored raccoons has probability 0.72 of developing the intended immune response. Let $X$ be the number that develop the response.
\begin{enumerate}
	\item State the probability distribution of $X$ and identify its parameters.
	\item Calculate $P(X=10)$.
	\item Calculate $P(X\geq11)$.
	\item Calculate $P(8\leq X\leq12)$.
	\item Find the mean and standard deviation of $X$ and interpret the mean.
	\item Identify two assumptions required for the binomial model to be appropriate.
\end{enumerate}

\questionheading{39}
An insurance claims office receives an average of 1.8 new claims during an eight-hour workday. Assume that claims arrive according to a Poisson process at a constant rate.
\begin{enumerate}
	\item Find the Poisson parameter for a three-hour period.
	\item Calculate the probability of receiving no claims during a three-hour period.
	\item Calculate the probability of receiving at least two claims during a three-hour period.
	\item Find the Poisson parameter for five complete workdays and calculate the probability of receiving between 6 and 12 claims, inclusive.
	\item Identify one realistic feature of claim arrivals that could violate the constant-rate assumption.
\end{enumerate}

\newpage

\questionheading{40}
A wildlife disease program monitors animals from two regions. Region A contains 40\% of the monitored population and has infection prevalence 18\%. Region B contains the remaining 60\% and has infection prevalence 6\%. A diagnostic test has sensitivity 92\% and specificity 96\% in both regions. Let $A$ denote selection from Region A, let $I$ denote infection and let $+$ denote a positive test result.
\begin{enumerate}
	\item Use the law of total probability to calculate $P(I)$.
	\item Calculate $P(+\mid A)$ and $P(+\mid A^c)$.
	\item Calculate the overall probability $P(+)$.
	\item Given a positive result, calculate $P(A\mid +)$.
	\item Given a positive result, calculate $P(I\mid +)$.
	\item If 20 animals are independently selected from this mixed population and tested, state a probability model for the number of positive results and calculate the probability of at least three positive results.
	\item Find the expected number of positive results among the 20 tests.
\end{enumerate}


\end{document}
